\section{Introduction}\label{sec:intro}

Deciding where to open distribution centers and which center serves each demand region is a
recurring strategic decision in logistics. When each region must be served by exactly one
center and centers have limited capacity, the decision is the single-source capacitated
facility location problem (SSCFLP). It is NP-hard, and the single-sourcing requirement makes it
markedly harder than its splittable counterpart: once the open set is fixed, the remaining
assignment is still a generalized assignment problem.

Two lines of work address large SSCFLP instances. Matheuristics restrict the problem and solve
the restriction exactly: kernel search ranks variables by the linear relaxation and solves a
sequence of restricted MILPs \citep{guastaroba2014}, and large neighborhood search (LNS)
repeatedly frees part of the solution and re-optimizes it \citep{kong2021}. More recently,
machine learning has been proposed to accelerate exact and heuristic solvers, either by
learning decisions inside branch-and-bound \citep{gasse2019,gupta2020}, by predicting which
variables to fix or explore \citep{nair2020,han2023}, by learning which neighborhood an LNS
should destroy \citep{song2020,sonnerat2021,huang2023}, or by learning which subproblem a local
search should delegate to an existing solver \citep{li2021delegate}.

Our starting point was the most direct of these ideas for facility location: let a graph neural
network rank candidate facilities and solve the problem restricted to the top-ranked ones. Two
observations changed the direction of the study. First, ranking quality did not translate into
safe pruning. On validation instances, the GNN ranked facilities well (AUC 0.95), yet keeping
the best known solution inside the restricted problem in 80\% of the instances required
retaining about 80\% of the candidates; the LP relaxation alone did almost as well, without
training (Section~\ref{sec:res-pruning}). Capacity forces slack to be kept, and under single
sourcing the slack is spread over many facilities. Second, decomposing the customers into
clusters and solving each cluster independently was fast but wrong: the union of the cluster
solutions violated capacity in every instance we examined and, after repair, was 19--29\% more
expensive than the reference.

The study is organized around four questions. RQ1: does progressive restriction of the
facilities improve the quality available over the budget, against the full solver and classical
controls? RQ2: does replacing only the classical ranking by a learned one improve the result?
RQ3: does adding the coordinated neighborhood search after the restriction improve the final
quality or the anytime behavior? RQ4 (exploratory): do repetition, neighborhood reach and
coordination explain part of the observed limitations? RQ1--RQ3 are answered by a pre-registered
closed test; RQ4 and everything about subproblem selection rest on validation pilots.

These observations redirected the study in two ways. First, we use the ranking to \emph{guide} a
search that never leaves the feasible region, rather than to remove parts of the problem.
Second, since the LP relaxation ranked facilities as safely as the GNN, we ask of every learned
component whether a learning-free counterpart with the same function does as well. We study two
mechanisms, each in a learned and a learning-free variant:
\begin{enumerate}
\item an \textbf{adaptive domain expansion} that solves restricted problems over 20, 40, 60 and
100\% of the ranked facilities and stops early when LP reduced costs certify that no excluded
facility can enter a solution better than the incumbent. The ranking is given by a GNN or by the
LP relaxation; kernel search \citep{guastaroba2014} is the classical baseline; and
\item a \textbf{cluster-oriented LNS (CLNS)} that re-optimizes one subproblem at a time (a
cluster, the boundary between two clusters, the customers of an expensive facility, or the
customers with the largest LP regret) with the rest of the solution fixed and the capacities
reduced by the fixed load. Every accepted move is feasible for the full problem and each fixed
cost is charged once. The subproblem is chosen by rotation, by adaptive operator weights, by a
learned regressor or by disagreement with the GNN, and a memory of failed subproblems measures
and avoids repetitions.
\end{enumerate}

\paragraph{Contributions.}
\begin{itemize}
\item \textbf{A controlled comparison} of learned and learning-free search-space reduction for the
SSCFLP, in which each learned component has a counterpart of the same function under the same
budget. A pre-registered test answers RQ1--RQ3 (Section~\ref{sec:res-test}); exploratory pilots
address RQ4 and subproblem selection (Sections~\ref{sec:res-pilot} and~\ref{sec:res-controls}).
\item \textbf{A characterization of failure modes} under capacity coupling: ranking-based pruning
that cannot be both small and safe, uncoordinated decomposition that produces infeasible unions,
and a neighborhood search in which more than half of the selected subproblems had already been
tried without success (Sections~\ref{sec:res-pruning}, \ref{sec:res-decomp}
and~\ref{sec:res-controls}).
\item \textbf{A reproducible implementation with delimited guarantees}: a coordinated
decomposition whose accepted moves are feasible for the full problem, a reduced-cost stopping
rule that lets a restricted solve prove optimality, and an evaluation protocol with pinned cores,
two clocks, frozen data and registered hypotheses (Sections~\ref{sec:method}
and~\ref{sec:experiments}). The reduced-cost argument and the idea of remembering unproductive
subproblems are classical; we use them, we do not claim them.
\end{itemize}
Appendix~\ref{app:reconstructions} reports, as supplementary material with its own scope, what we
observed when reconstructing five recent learning-for-optimization approaches.

In the pre-registered closed test, run twice on different hardware, restricting the search
space by the LP ranking about halved the primal integral of the solver on the full model, while
the same mechanism driven by the learned ranking had a significantly worse primal integral than
its LP-driven counterpart in both rounds and, on dedicated hardware, was numerically above the full
solver. On real instances with up to
850 customers, only the methods with the coordinated neighborhood search came close to the best
known solution on the two largest.
Section~\ref{sec:related} reviews related work; Section~\ref{sec:problem} states the problem, its
relaxation and the capacity coupling; Section~\ref{sec:method} describes the mechanisms;
Section~\ref{sec:experiments} gives the experimental design; Section~\ref{sec:results} reports
the results; Section~\ref{sec:discussion} discusses them and Section~\ref{sec:conclusion}
concludes.
